\scanpage{45}
D'après le théor. d'Artin--Schreier--Witt, on a
\[
H^1(A,\mathbb Z/p)\simeq H^1(A,\underline O_A)^{\unsure{(\varphi)|(p)}}_{\deleted{\unsure{\varphi}}}
\]
éléments de $H^1(A,\underline O_A)$ tels que
\[Fx=x\]
\[
H^1(A,\mathbb Z/p^n)\simeq H^1(\unsure{A},\underline O_A(W_n))^{\unsure{(\varphi)|(p)}}
\]
--- $H^1(A,\underline O_A(W_n))$ tels que
\[Fx=x\]
\[
\boxed{\operatorname{Hom}(\pi_1^p(A),\mathbb Z_p)\simeq H^1(A,\deleted{\unsure{\mathbb Q_p}}\mathbb Z_p)
\simeq H^1(A,\unsure{\mathbb W|W})^{(\varphi)}_{\deleted{\unsure{\varphi}}}}
\]
\inserted{d'où} \deleted{\unsure{on a}} on déduit un \deleted{accouplement} isomorphisme
\[
\boxed{\deleted{\pi_1^p(A)\times H^1(A,\mathbb W)_s\longrightarrow\mathbb Z_p}}
\]
\[
\boxed{\deleted{\operatorname{Hom}_{\mathrm{cont}}(\pi_1^p(A),\Lambda)\simeq H^1(A,\mathbb W)_s}}
\]
qui est une dualité \deleted{\unsure{de}} de $\mathbb Z_p$-modules libres de type fini

D'autre part, utilisant encore le fait que tout revêtement de $A$ est une variété abélienne, on trouve aussi
\[\boxed{\pi_1^p(A)\simeq T_p(A)}\]
d'où en les composant
\[\boxed{T_p(A)\times H^1(A,\mathbb W)^{(\varphi)}\longrightarrow\mathbb Z_p}\]
qui est une dualité de $\mathbb Z_p$-modules libres de type fini.

\deleted{qu'on peut aussi écrire} d'où on déduit
\[
\boxed{\operatorname{Hom}_{\mathrm{cont}}(T_p(A),\deleted{\unsure{\Lambda}}W_\infty(k))\simeq H^1(A,\mathbb W)_s}
\]
\margin{44}
